\section*{Bab \ref{ch:b2:acyl-substitution} --- Turunan Asam Karboksilat}

\begin{solution}{exo:b2:acyl-substitution:1}
Etanamida $<$ etil etanoat $<$ anhidrida etanoat $<$ etanoil klorida.
\end{solution}

\begin{solution}{exo:b2:acyl-substitution:2}
Asam etanoat (+ \ce{HCl}); etil etanoat; etanamida (+ amonium klorida,
dengan ekuivalen kedua amonia yang menangkap \ce{HCl}); anhidrida etanoat
(+ \ce{NaCl}); \textit{N},\textit{N}-dimetiletanamida (+ trietilamonium
klorida).
\end{solution}

\begin{solution}{exo:b2:acyl-substitution:3}
\ce{C6H5COOCH3 + NaOH -> C6H5COONa + CH3OH}. $13.6/136.15 = \qty{0.0999}{mol}$, \ce{NaOH}
sebanyak itu pula: \qty{4.00}{g}.
\end{solution}

\begin{solution}{exo:b2:acyl-substitution:4}
$\gamma$-Butirolakton (oksolan-2-on): cincin beranggota lima yang terdiri atas empat karbon dan satu oksigen.
\end{solution}

\begin{solution}{exo:b2:acyl-substitution:5}
Amonia beradisi pada karbon karbonil; sebuah proton berpindah dari nitrogen
ke oksigen (atau ke gugus pergi); etoksida (sebagai etanol, setelah
transfer proton) pergi dan \ce{C=O} terbentuk kembali: etanamida dan etanol.
\end{solution}

\begin{solution}{exo:b2:acyl-substitution:6}
\ce{OH} fenolik terasilasi (katalisis asam): asam 2-asetoksibenzoat; produk sampingnya asam etanoat.
\end{solution}

\begin{solution}{exo:b2:acyl-substitution:7}
2-Metilpropan-2-ol. Adisi pertama menghasilkan propanon, yang lebih
elektrofilik daripada ester, dan langsung mengikat metil kedua.
\end{solution}

\begin{solution}{exo:b2:acyl-substitution:8}
Butana-1,4-diol: ester direduksi menjadi dua alkohol primer, dan cincinnya terbuka.
\end{solution}

\begin{solution}{exo:b2:acyl-substitution:9}
Satu ekuivalen: $x^2/(1 - x)^2 = 4$, $x = 2/3$. Tiga ekuivalen: $x^2/[(1 - x)(3 - x)] = 4$, $3x^2 - 16x
+ 12 = 0$, $x = 0.90$. Mengeluarkan air begitu air terbentuk (Dean--Stark)
mendorong reaksi lebih jauh.
\end{solution}

\begin{solution}{exo:b2:acyl-substitution:10}
\ce{SOCl2} menghasilkan benzoil klorida; dengan dua ekuivalen dietilamina
(atau satu ditambah trietilamina) senyawa itu menghasilkan amida. Asam dan
amina yang dicampur langsung mengalami reaksi asam--basa: dietilamonium
benzoat, yang hanya menghasilkan amida pada pemanasan kuat.
\end{solution}

\begin{solution}{exo:b2:acyl-substitution:11}
Tiga KOH per triolein: $3 \times 56.11/885.4 = \qty{0.190}{g}$ per gram, yaitu
bilangan penyabunan 190.
\end{solution}

\begin{solution}{exo:b2:acyl-substitution:12}
\ce{NaBH4}: hanya keton, menjadi alkohol sekunder. \ce{LiAlH4}: keton menjadi
alkohol sekunder, ester menjadi alkohol primer (dan melepaskan bagian
alkoholnya), amida menjadi amina.
\end{solution}

\begin{solution}{pb:b2:acyl-substitution:1}
\textbf{1.} Asam oleat \ce{C18H34O2}; gliserol \ce{C3H8O3}.
\textbf{2.} \ce{C57H104O6}: $57 \times 12.011 + 104 \times 1.008 + 6 \times 15.999 = \qty{885.45}{g/mol}$.
\textbf{3.} Tiga.
\textbf{4.} Ikatan rangkap cis membengkokkan rantai-rantainya, yang tersusun
kurang rapat: gaya antarmolekul lemah, titik leleh rendah.
\textbf{5.} Gliserol dan tiga natrium oleat, yaitu sabun.
\textbf{6.} Campuran metil ester asam-asam lemak.
\textbf{7.} \ce{C57H104O6 + 3CH3OH -> C3H8O3 + 3C19H36O2}.
\textbf{8.} Ion metoksida, yang terbentuk dari metanol dan basa, adalah
nukleofil yang jauh lebih kuat daripada metanol; ion ini menyerang karbonil
ester.
\textbf{9.} Karbon karbonil salah satu gugus ester yang membawa \ce{O-},
\ce{OCH3} dari metoksida, dan oksigen gliseril: tetrahedral.
\textbf{10.} \omterm{def:b2:acyl-substitution:transesterification}{Transesterifikasi} adalah kesetimbangan dengan tetapan sekitar 1;
metanol berlebih menggesernya ke arah metil ester.
\textbf{11.} Gliserol memisah sebagai lapisan bawah: mengeluarkan sebuah
produk mendorong kesetimbangan.
\textbf{12.} Tidak: metoksida diregenerasi pada setiap pertukaran.
\textbf{13.} Asam itu menetralkan basa dan menghasilkan natrium oleat
(sabun): katalisnya terpakai.
\textbf{14.} Air menyabunkannya: hidroksida (dari air dan metoksida)
memotong ester menjadi sabun dan metanol, secara ireversibel.
\textbf{15.} Sabun itu menghabiskan basa, mengemulsikan campuran, dan
menghalangi lapisan gliserol untuk memisah.
\textbf{16.} Dengan \omterm{def:b2:acyl-substitution:esterification}{esterifikasi} terkatalisis asam bersama metanol lebih
dahulu, yang mengubah asam-asam bebas menjadi metil ester.
\textbf{17.} Air menghidrolisis ester dan mengubah katalis menjadi
hidroksida, yang menyabunkan.
\textbf{18.} $10^6/885.45 = \qty{1129}{mol}$.
\textbf{19.} $3 \times 1129 \times 32.04 = \qty{108.6}{kg}$.
\textbf{20.} $3 \times 1129 \times 296.50 = \qty{1004.6}{kg}$.
\textbf{21.} $1000 + 108.6 = \qty{1108.6}{kg}$ masuk; $104.0 + 1004.6 = \qty{1108.6}{kg}$ keluar.
\textbf{22.} \qty{20}{kg} asam oleat, $20\,000/282.47 = \qty{70.8}{mol}$, \ce{NaOH} sebanyak itu pula:
\qty{2.83}{kg}.
\textbf{23.} $1129 \times 92.09 = \boldsymbol{\approx \qty{104}{kg}}$ gliserol per ton.
\end{solution}
